\chapter{ระบบความปลอดภัย Zero-Trust และการดีพลอยระดับองค์กร}

ความปลอดภัยของข้อมูลบนระบบคลาวด์ ต้องยึดหลักการไม่ไว้วางใจอินพุตใดๆ จากฝั่งไคลเอนต์ (Zero-Trust Security) และตรวจสอบสิทธิ์ในทุกคำขออย่างเคร่งครัด

\section{การออกแบบ Security Rules เชิงลึก}

การเขียนกฎความปลอดภัยใน Cloud Firestore และ Storage ต้องตรวจสอบทั้งการยืนยันตัวตน (Authentication) และการอนุญาตเข้าถึงสิทธิ์ (Authorization):

\begin{lstlisting}[language=bash, caption={Example Firestore Security Rule}]
match /experiments/{docId} {
  allow read: if request.auth != null;
  allow write: if request.auth != null && 
                  request.auth.uid == resource.data.authorId;
}
\end{lstlisting}

\section{สถาปัตยกรรม Firebase App Hosting}

สำหรับเว็บแอปพลิเคชันยุคใหม่ที่สร้างด้วย Next.js หรือ Angular สถาปัตยกรรม App Hosting จะทำการจัดการการเรนเดอร์ฝั่งเซิร์ฟเวอร์ (Server-Side Rendering: SSR) และสร้าง CDN แคชทั่วโลกโดยอัตโนมัติเมื่อมีการผลักโค้ด (Push) ขึ้นสู่ GitHub

\section{การตรวจสอบความเสถียรด้วย Crashlytics}

การผสาน Firebase Crashlytics ช่วยให้ผู้ดูแลระบบมองเห็นปัญหาแอปพลิเคชันหยุดทำงาน (Crash-Free Users Percentage) พร้อมสแต็กเทรซ (Stack Trace) เชิงลึก ทำให้แก้ไขจุดบกพร่องได้ทันทีก่อนส่งผลกระทบต่อผู้ใช้งานในวงกว้าง
